% !TeX encoding = UTF-8
% Build from repository root:
% latexmk -pdfdvi -e '$latex=q/uplatex %O -interaction=nonstopmode -halt-on-error %S/' -e '$dvipdf=q/dvipdfmx %O -o %D %S/' -outdir=paper/shared/build paper/shared/amber_notation.tex
\documentclass[uplatex,dvipdfmx,fontsize=11pt,baselineskip=17pt,paper=a4]{jlreq}
% Thesis typography and existing indicator macro; no shared preamble exists.
\usepackage{mathtools,amssymb,booktabs,array}
\usepackage{newtxtext,newtxmath}
\usepackage[haranoaji]{pxchfon}
\usepackage[top=20mm,bottom=20mm,left=18mm,right=18mm]{geometry}
\newcommand{\ind}{\mathbf{1}}
\newcommand{\AC}{\mathrm{AC@}}
\newcommand{\Avg}{\mathrm{Avg@}}
\newcommand{\head}[1]{\par\bigskip\noindent{\large\bfseries #1}\par\smallskip}
\setlength{\parindent}{0pt}
\setlength{\parskip}{5pt}
\setlength{\abovedisplayskip}{10pt}
\setlength{\belowdisplayskip}{10pt}
\renewcommand{\arraystretch}{1.3}
\begin{document}
\begin{center}
{\LARGE\bfseries AMBER 記号と計算の補足資料}\\[3pt]
入力から評価まで、具体例で読む記法ガイド
\end{center}
修論本文とは独立して読める補足資料。本文・図・スライドで共通に用いる記法を、計算の順に整理する。

以下の数値例は説明用の仮想例であり、実験結果ではない。

上付き $(N)$ は正常期間、$(A)$ は障害期間、
$(\mathrm{cf})$ は反実仮想（counterfactual）を表すラベルであり、累乗ではない。

\head{1. 添字と入力系列}
\begin{tabular}{@{}p{48mm}p{122mm}@{}}
\toprule
記号 & 意味 \\
\midrule
$i=1,\ldots,M$ ／ $s=1,\ldots,S$ & メトリクスの添字・総数 ／ サービスの添字・総数 \\
$t=1,\ldots,T$ ／ $t_F$ & 観測時刻（離散的な標本番号）／ 障害発生時刻 \\
$T_N=t_F-1$ & 正常期間長。$t_F$ の直前までを正常期間とする。\\
$T_A=T-t_F+1$ & 障害期間長。$t_F$ を含み、$T=T_N+T_A$。\\
$h=1,\ldots,T_A$ & 障害期間の局所時刻（予測の先行ステップ数）。$t=t_F+h-1$。\\
$P,\ c_0,\ \phi_p$ & AR次数、切片、ラグ $p$ の係数。正常期間だけで推定。\\
\bottomrule
\end{tabular}
\[
\mathbf{x}_i=(x_{i,1},\ldots,x_{i,T})^\top
\longmapsto
\begin{cases}
\mathbf{x}_i^{(N)}=(x_{i,1}^{(N)},\ldots,x_{i,T_N}^{(N)})^\top,&x_{i,t}^{(N)}=x_{i,t},\\
\mathbf{x}_i^{(A)}=(x_{i,1}^{(A)},\ldots,x_{i,T_A}^{(A)})^\top,&x_{i,h}^{(A)}=x_{i,t_F+h-1}.
\end{cases}
\]
以下、1メトリクスの計算では $i$ を省略する。切片は実験ケースの添字 $c$ と区別して $c_0$ とする。
窓を切り出す場合、$t$ は採用した観測窓内で付け直した番号とする。

\textbf{例：観測時刻と障害期間の添字}

$T=10$、$t_F=7$ なら、$T_N=6$、$T_A=4$ となる。
正常系列は $(x_{i,1},\ldots,x_{i,6})^\top$、障害系列は
$(x_{i,7},\ldots,x_{i,10})^\top$ である。
障害期間の $h=1$ は元の $t=7$、$h=4$ は $t=10$ に対応する。

\newpage
\head{2. 正常期間だけに基づく入力の標準化}
\[
\mu_x=\frac{1}{T_N}\sum_{t=1}^{T_N}x_t^{(N)},\qquad
\sigma_x=\sqrt{\frac{1}{T_N}\sum_{t=1}^{T_N}(x_t^{(N)}-\mu_x)^2}.
\]
\[
 y_t^{(N)}=\frac{x_t^{(N)}-\mu_x}{\sigma_x},\quad
 y_h^{(A)}=\frac{x_h^{(A)}-\mu_x}{\sigma_x}.
\]
$x$ は観測値、$y$ は標準化後の観測値。両期間に同じ正常時の平均・尺度を使う。
上式は $\sigma_x>0$ の通常ケースを示す（実装では定数系列等への数値的保護を行う）。

\textbf{例：同じ変換を両期間に使う}

正常期間で $\mu_x=10$、$\sigma_x=2$ と求まったとする。
正常値 $x_t^{(N)}=12$ は $y_t^{(N)}=1$、障害値 $x_h^{(A)}=16$ は $y_h^{(A)}=3$ となる。
障害期間だけで平均・標準偏差を求め直すと、検出したい水準変化まで消してしまう場合がある。

\head{3. 正常期間の1期先予測と反実仮想予測}
\[
\hat y_{t\mid t-1}^{(N)}=c_0+\sum_{p=1}^{P}\phi_p y_{t-p}^{(N)},
\]
\[
r_t^{(N)}=y_t^{(N)}-\hat y_{t\mid t-1}^{(N)},\quad t=P+1,\ldots,T_N.
\]
\textbf{1期先予測}は、$t-1$ 以前の実測値をラグに使って時刻 $t$ を予測すること。
$r_t^{(N)}$ は正常期間の1期先予測残差である。
\[
\hat y_h^{(\mathrm{cf})}=c_0+\sum_{p=1}^{P}\phi_p\tilde y_{h-p}^{(\mathrm{cf})},
\]
\[
\tilde y_j^{(\mathrm{cf})}=\begin{cases}
y_{T_N+j}^{(N)},&1-P\le j\le0,\\
\hat y_j^{(\mathrm{cf})},&j\ge1,
\end{cases}
\]
\[
r_h^{(A)}=y_h^{(A)}-\hat y_h^{(\mathrm{cf})}.
\]
$\hat y_h^{(\mathrm{cf})}$ は正常状態の継続を仮定した予測、$r_h^{(A)}$ は障害期間の反実仮想予測誤差。
正常期間末尾を起点に再帰予測し、\textbf{障害期間の実測値はラグに戻さない}。

\textbf{例：1期先予測と再帰予測の違い}

$P=1$、$c_0=0$、$\phi_1=0.5$、正常期間末尾 $y_{T_N}^{(N)}=2$ とする。
反実仮想予測は $\hat y_1^{(\mathrm{cf})}=1$、$\hat y_2^{(\mathrm{cf})}=0.5$。
実測が $y_1^{(A)}=4$、$y_2^{(A)}=3$ なら、誤差はそれぞれ $3$、$2.5$ となる。
2期目の予測には実測の $4$ ではなく、前の予測値 $1$ を使う。

\newpage
\head{4. 予測不確実性と誤差の標準化}
\[
\psi_0=1,\qquad \psi_j=\sum_{p=1}^{\min(P,j)}\phi_p\psi_{j-p}\ (j\ge1),\qquad
q_h=\sqrt{\sum_{j=0}^{h-1}\psi_j^2},\qquad q_1=1.
\]
$q_h$ は予測先 $h$ に応じた不確実性倍率（horizon uncertainty）。$\psi_j$ はARのインパルス応答。
正常残差の中央値 $a=\operatorname{median}_t r_t^{(N)}$ と正常残差から定めた尺度 $b>0$ を使い、
\[
\boxed{z_t^{(N)}=\frac{r_t^{(N)}-a}{b}},\qquad
\boxed{z_h^{(A)}=\frac{r_h^{(A)}-a}{bq_h}}.
\]
$b$ はMADを主とし標準偏差・尺度下限で保護した正常時の尺度。
ここでは各予測先の分散で補正する設定を記す。$z$ は標準化誤差であり、入力の $y$ と区別する。

\textbf{例：遠い予測先の誤差を割り引く}

前ページのAR(1)では $\psi_j=0.5^j$ なので、
\[
q_1=1,\qquad q_2=\sqrt{1+0.5^2}=\sqrt{1.25}\simeq1.118.
\]
$a=0$、$b=1$ とすると、$r_1^{(A)}=3$ は $z_1^{(A)}=3$、
$r_2^{(A)}=2.5$ は $z_2^{(A)}=2.5/1.118\simeq2.236$ となる。
予測先が遠いほど増えるモデル上の不確実性を、異常の大きさと区別する補正である。

\head{似た記号の読み分け}
\begin{tabular}{@{}p{35mm}p{130mm}@{}}
\toprule
記号 & 役割 \\
\midrule
$x$ & 元の単位を持つ観測値。\\
$y$ & 正常時の平均・尺度で標準化した、ARへの入力。\\
$r$ & 観測と予測の差。予測がどれだけ外れたか。\\
$z$ & 正常時の残差尺度と予測不確実性で標準化した誤差。\\
\bottomrule
\end{tabular}

入力の標準化に使う $\sigma_x$ と、残差の標準化に使う $b$ は別の尺度である。
また、$q_h$ は平均や分散を新たに障害データから学習する量ではなく、正常期間で推定したAR係数から計算する。

\newpage
\head{5. 仮説比較とメトリクス証拠スコア}
\[
D_N=\{z_t^{(N)}:t=P+1,\ldots,T_N\},\qquad
D_A=\{z_h^{(A)}:h=1,\ldots,T_A\}.
\]
\begin{tabular}{@{}p{30mm}p{140mm}@{}}
\toprule
記号 & 意味 \\
\midrule
$H_0$ & 両期間が平均・分散を共有する同一の正規分布に従うモデル。\\
$H_1$ & 両期間で別々の平均・分散を持つ正規分布を許すモデル。\\
$\mathrm{BF}_i$ & メトリクス $i$ における $H_0$ に対する $H_1$ のベイズ因子。\\
$E_i=\log\mathrm{BF}_i$ & メトリクス証拠スコア。大きいほど分布変化を支持する。\\
\bottomrule
\end{tabular}
\[
\mathrm{BF}_i=\frac{p(D_{N,i},D_{A,i}\mid H_1)}{p(D_{N,i},D_{A,i}\mid H_0)},\]
\[
E_i=\log p(D_{N,i},D_{A,i}\mid H_1)-\log p(D_{N,i},D_{A,i}\mid H_0).
\]
$p(\cdot\mid H)$ はパラメータを事前分布で積分した周辺尤度、$\log$ は自然対数。
$E_i$ は原因である確率ではなく、正常時からの変化を示す証拠の強さである。

\head{6. サービス集約と順位}
サービス $s$ に属するメトリクス添字集合を $\mathcal M_s$、その要素数を $n_s=|\mathcal M_s|$ とする。
その証拠スコアを $E_{s,(1)}\ge\cdots\ge E_{s,(n_s)}$ と降順に並べる。
\[
K_s=\min(3,n_s),\qquad
\boxed{R_s=\frac{1}{K_s}\sum_{j=1}^{K_s}E_{s,(j)}}\quad(n_s\ge1).
\]
\[
\boldsymbol{\pi}=(\pi_1,\ldots,\pi_S),\quad R_{\pi_1}\ge\cdots\ge R_{\pi_S}.
\]
$R_s$ はサービスRCAスコア（上位3件の平均、\texttt{mean\_top3}）。
$\pi_k$ は第 $k$ 位のサービス。$S$ は順位付け対象のサービス数とし、同点は実装の規則で順序を確定する。
$K_s$ は集約件数、以下の $K$ は評価する順位の上限であり、別の記号である。

\textbf{例：メトリクスからサービスへ}

サービスAの証拠スコアが $8,5,2,-1$ なら、$R_{\mathrm A}=(8+5+2)/3=5$。
サービスBが $6,4$ の2メトリクスだけなら、$R_{\mathrm B}=(6+4)/2=5$。
3件未満のサービスは、利用できる件数で平均する。
サービスCのスコアが $7$ ならCが第1位となり、AとBは同点規則で順序を決める。

\newpage
\head{7. 正解順位とケース単位の評価}
\begin{tabular}{@{}p{40mm}p{130mm}@{}}
\toprule
記号 & 意味 \\
\midrule
$c=1,\ldots,C$ & 実験ケースの添字と総数。\\
$g_c$ ／ $\boldsymbol{\pi}_c$ & ケース $c$ の正解サービス（ground truth）／ 出力された順位列。\\
$r_c=\min\{k:\pi_{c,k}=g_c\}$ & 正解サービスの順位。順位列に含まれなければ $r_c=\infty$。\\
\bottomrule
\end{tabular}
\[
\boxed{\AC K(c)=\ind[r_c\le K]},
\]
\[
\boxed{\Avg K(c)=\frac{1}{K}\sum_{k=1}^{K}\AC k(c)}
=\begin{cases}(K-r_c+1)/K,&r_c\le K,\\0,&r_c>K.\end{cases}
\]
$K$ は正の整数。AC@Kは正解が上位 $K$ 件に入るかを0/1で表し、Avg@Kは上位内での順位も反映する。
例えば $r_c=3$ なら AC@5 $=1$、Avg@5 $=3/5=0.6$。
ここでは単一正解を仮定する。複数正解の場合は、正解集合内の最上位の順位を $r_c$ とする。

\textbf{例：AC@5とAvg@5の違い}

\begin{center}
\begin{tabular}{ccc}
\toprule
正解順位 $r_c$ & AC@5 & Avg@5 \\
\midrule
1 & 1 & 1.0 \\
2 & 1 & 0.8 \\
3 & 1 & 0.6 \\
4 & 1 & 0.4 \\
5 & 1 & 0.2 \\
6以上・圏外 & 0 & 0.0 \\
\bottomrule
\end{tabular}
\end{center}

正解が第3位なら、AC@1からAC@5までは $(0,0,1,1,1)$。
この5個を平均するのでAvg@5は $0.6$ となる。
AC@5は第1位と第5位を同じ成功として数えるが、Avg@5はより上位の正解を高く評価する。

\newpage
\head{8. ケース平均とデータセット間のマクロ平均}
\[
\AC K=\frac{1}{C}\sum_{c=1}^{C}\AC K(c),\qquad
\Avg K=\frac{1}{C}\sum_{c=1}^{C}\Avg K(c)=\frac{1}{K}\sum_{k=1}^{K}\AC k.
\]
データセット $d=1,\ldots,D$ のケース集合を $\mathcal C_d$、件数を $C_d=|\mathcal C_d|>0$ とする。
\[
\AC K_d=\frac{1}{C_d}\sum_{c\in\mathcal C_d}\AC K(c),\qquad
\Avg K_d=\frac{1}{C_d}\sum_{c\in\mathcal C_d}\Avg K(c),
\]
\[
\boxed{\mathrm{macro\text{-}AC@}K=\frac{1}{D}\sum_{d=1}^{D}\AC K_d},\qquad
\boxed{\mathrm{macro\text{-}Avg@}K=\frac{1}{D}\sum_{d=1}^{D}\Avg K_d}.
\]
マクロ平均はデータセットを等重みで平均する。全ケース平均とは重みが異なり、各 $C_d$ が等しい場合は一致する。
\textbf{$E_i\to R_s\to\boldsymbol{\pi}$ は推論の出力、AC@K・Avg@Kは正解と比較する評価指標}である。
\textbf{例：全ケース平均とマクロ平均}

2つのデータセットで、AC@1の成功件数が次のようになったとする。
\begin{center}
\begin{tabular}{lrrr}
\toprule
データセット & ケース数 & 成功件数 & AC@1 \\
\midrule
A & 100 & 90 & 0.9 \\
B & 20 & 10 & 0.5 \\
\bottomrule
\end{tabular}
\end{center}
\[
\text{全ケース平均}=\frac{90+10}{100+20}\simeq0.833.
\]
\vspace{6pt}
\[
\text{macro-AC@1}=\frac{0.9+0.5}{2}=0.7.
\]
全ケース平均ではケース数の多いAが大きな重みを持つ。
マクロ平均ではAとBを同じ重みで扱う。Avg@Kにも同じ平均方法を適用する。

\head{読み返すときの順序}
観測 $x$ を正常・障害期間に分け、正常時の情報で $y$ に変換する。
AR予測との差 $r$ を尺度補正して $z$ とし、期間間の分布変化を $E_i$ で測る。
その後、サービススコア $R_s$ に集約して順位 $\boldsymbol{\pi}$ を出す。
最後に正解 $g_c$ と比較し、AC@K・Avg@Kを計算する。

\end{document}
